\documentclass[final,hyperref={pdfpagelabels=false}]{beamer}
\usepackage[orientation=landscape,size=custom,width=121.92,height=91.44,scale=1.4]{beamerposter}
\usetheme{default}
\usepackage{amsmath,amssymb,amsthm}
\usepackage{mathtools}
\usepackage{bm}
\usepackage{tikz}
\usetikzlibrary{arrows,arrows.meta,positioning,shapes,calc,decorations.pathreplacing}
\usepackage{graphicx}
\usepackage{booktabs}
\usepackage{algorithm}
\usepackage{algorithmic}

\graphicspath{{../../figures/}}

\definecolor{harvardcrimson}{RGB}{165,28,48}
\definecolor{softblue}{RGB}{65,105,225}
\definecolor{softgreen}{RGB}{34,139,34}
\definecolor{softorange}{RGB}{255,140,0}
\definecolor{lightgray}{RGB}{245,245,245}
\definecolor{darkgray}{RGB}{60,60,60}
\definecolor{lightorange}{RGB}{255,200,120}

\setbeamercolor{headline}{bg=harvardcrimson,fg=white}
\setbeamercolor{block title}{bg=harvardcrimson,fg=white}
\setbeamercolor{block body}{bg=lightgray,fg=darkgray}
\setbeamercolor{block alerted title}{bg=softorange,fg=white}
\setbeamercolor{block alerted body}{bg=lightorange!10,fg=darkgray}
\setbeamerfont{block title}{size=\Large,series=\bfseries}
\setbeamerfont{headline title}{size=\Huge,series=\bfseries}
\setbeamertemplate{navigation symbols}{}

\newlength{\sepwidth}
\newlength{\colwidth}
\setlength{\sepwidth}{0.02\paperwidth}
\setlength{\colwidth}{0.30\paperwidth}

\newcommand{\E}{\mathbb{E}}
\newcommand{\Var}{\mathrm{Var}}

\begin{document}
\begin{frame}[t]

\begin{beamercolorbox}[wd=\paperwidth,ht=10cm]{headline}
\vspace{1cm}
\centering
{\fontsize{72}{80}\selectfont\bfseries MALA-Guided Mirror Descent}\\[0.8cm]
{\fontsize{42}{50}\selectfont Inference-Time Tilting of Diffusion Policies for Online RL}\\[1cm]
{\fontsize{36}{44}\selectfont Philp Nielsen}\\[0.5cm]
\end{beamercolorbox}

\vspace{1cm}

\begin{columns}[t]

\begin{column}{\colwidth}

\begin{block}{\Large Mirror Descent as an Inference-Time Policy Update}
\small
Mirror descent improves a policy by exponentially tilting it toward high-value actions:
\[
\pi'(a\mid s) \propto \pi_\theta(a\mid s)\,\exp\!\bigl(\eta\,Q(s,a)\bigr).
\]

\vspace{0.3em}
MGMD uses a diffusion policy as the base distribution $\pi_\theta$ and performs this improvement step at sampling time.

\vspace{0.3em}
\textbf{Why diffusion policies are a good fit:}
\begin{itemize}
    \item \textbf{Expressive:} they represent multimodal action distributions well.
    \item \textbf{Trainable:} the policy is still learned by standard denoising score matching.
    \item \textbf{Deployable:} all critic information is injected at inference time, then distilled back into the policy through replay.
\end{itemize}

\vspace{0.3em}
The practical challenge is not writing down $\pi'$, but sampling from it accurately during reverse diffusion.
\end{block}

\vspace{0.8cm}

\begin{block}{\Large Diffusion Policies, Tweedie Reconstruction, and Guidance}
\small
A diffusion policy starts from Gaussian noise and denoises toward an action:
\[
q(a_t\mid a_0)=\mathcal{N}\!\bigl(\sqrt{\bar\alpha_t}\,a_0,\,(1-\bar\alpha_t)I\bigr),
\qquad
a_T \to a_{T-1} \to \cdots \to a_0.
\]

\vspace{0.3em}
The network predicts noise, which gives the Tweedie clean-action estimate
\[
\hat a_0(a_t)=\frac{a_t-\sqrt{1-\bar\alpha_t}\,\hat\epsilon_\theta(a_t,s,t)}{\sqrt{\bar\alpha_t}}.
\]

\vspace{0.3em}
A natural training-free guidance rule evaluates the critic only on clean actions:
\[
\hat\epsilon_t^{\mathrm{Tw}}(a_t)
=
\hat\epsilon_\theta(a_t,s,t)
-
\eta\sqrt{1-\bar\alpha_t}\,\nabla_{a_t}Q\!\bigl(s,\hat a_0(a_t)\bigr).
\]

\vspace{0.3em}
This is attractive because the critic never has to model noisy actions, but it does \emph{not} give the exact score of the tilted noisy marginals for intermediate $t$.
\end{block}

\vspace{0.8cm}

\begin{block}{\Large Why MGMD Needs MCMC Corrections}
\small
The exact reverse process for the tilted policy would require the noisy tilted marginals $\tilde p_t$, not just the clean tilt at $t=0$.
Instead, MGMD treats each noise level as a separate energy-based target:
\[
p_t(a_t\mid s) \propto \exp\!\bigl(-E_t^{\mathrm{guided}}(a_t,s)\bigr),
\qquad
E_t^{\mathrm{guided}}(a_t,s)=E_\theta(s,a_t,t)-\eta\,Q\!\bigl(s,\hat a_0(a_t)\bigr).
\]

\vspace{0.3em}
This viewpoint gives the right endpoints:
\begin{itemize}
    \item at high noise, the target is close to the Gaussian base distribution,
    \item as $t \to 0$, the target approaches $\pi'(a\mid s) \propto \pi_\theta(a\mid s)e^{\eta Q(s,a)}$.
\end{itemize}

\vspace{0.3em}
MGMD therefore corrects the sampler \textbf{at every noise level} with MALA, rather than relying on plain guidance alone.
\end{block}

\vspace{0.8cm}

\begin{block}{\Large Adaptive Guidance Strength}
\small
The current project also adapts the tilt magnitude instead of fixing it forever. A simple scale-invariant rule is the KL-budget update
\[
\eta_k = \sqrt{\frac{2\Delta}{\E_{s\sim d_{\pi_k}}\!\left[\Var_{a\sim\pi_k}\bigl(A_k(s,a)\bigr)\right]}}.
\]

\vspace{0.2em}
In the code, this can be combined with a one-step distribution-shift correction that tracks advantage moments with EMAs. The effect is intuitive: guidance becomes weaker when critic uncertainty or policy-shift risk is high, and stronger when the critic signal is consistently informative.
\end{block}

\end{column}

\begin{column}{\colwidth}

\begin{block}{\Large Energy-Based Parameterization}
\small
Metropolis--Hastings needs \textbf{energy differences}, not just score vectors. MGMD therefore uses an energy-based diffusion policy:
\[
\hat\epsilon_\theta(s,a_t,t)
\coloneqq
\sqrt{1-\bar\alpha_t}\,\nabla_{a_t}E_\theta(s,a_t,t).
\]

\vspace{0.3em}
The training objective is still the standard diffusion loss:
\[
\mathcal{L}_{\mathrm{diff}}(\theta)
=
\E_{a_0,\epsilon,t}
\left[
\left\lVert
\hat\epsilon_\theta\!\left(s,\sqrt{\bar\alpha_t}a_0+\sqrt{1-\bar\alpha_t}\,\epsilon,t\right)-\epsilon
\right\rVert^2
\right].
\]

\vspace{0.3em}
This keeps diffusion training simple, while making both of the ingredients for MALA available:
\begin{itemize}
    \item gradients $\nabla_{a_t}E_t^{\mathrm{guided}}$ for Langevin proposals,
    \item energy differences $E_t^{\mathrm{guided}}(a)-E_t^{\mathrm{guided}}(a')$ for MH acceptance.
\end{itemize}
\end{block}

\vspace{0.8cm}

\begin{block}{\Large MALA Correction at Each Diffusion Level}
\small
At each noise level $t$, MGMD runs a short MALA chain targeting
\[
p_t(a_t\mid s) \propto \exp\!\bigl(-E_t^{\mathrm{guided}}(a_t,s)\bigr).
\]

\vspace{0.3em}
\textbf{Langevin proposal:}
\[
a_t' = a_t - s_t\,\nabla_{a_t}E_t^{\mathrm{guided}}(a_t,s) + \sqrt{2s_t}\,z,
\qquad z\sim\mathcal{N}(0,I).
\]

\textbf{MH correction:} with $\mu(a)=a-s_t\nabla E_t^{\mathrm{guided}}(a)$,
\[
\alpha(a,a')
=
\min\!\left\{1,
\exp\!\left(
E_t^{\mathrm{guided}}(a)-E_t^{\mathrm{guided}}(a')
+
\frac{\|a'-\mu(a)\|^2-\|a-\mu(a')\|^2}{4s_t}
\right)
\right\}.
\]

\vspace{0.3em}
Each diffusion level maintains its own step size $s_t$, adapted online toward the classical MALA target acceptance rate of \textbf{0.574}. This matters because the energy landscape is smooth at high noise and sharp near $t=0$.
\end{block}

\vspace{0.8cm}

\begin{block}{\Large Sampler: MALA Corrector + DDIM Predictor}
 \small
 \textbf{Algorithm 1:} tilted-policy sampler
 
 {\footnotesize
 \begin{algorithmic}[1]
 \STATE Sample $a_T \sim \mathcal{N}(0,I)$
 \FOR{$t=T,T-1,\ldots,1$}
     \STATE Form $E_t^{\mathrm{guided}}(a_t,s)$ using $\hat a_0(a_t)$
     \FOR{$m=1,\ldots,M$}
         \STATE Propose $a_t' \sim \mathcal{N}(a_t-s_t\nabla E_t^{\mathrm{guided}},2s_tI)$
         \STATE Accept or reject with the MH rule above
         \STATE Adapt $s_t$ toward acceptance rate $0.574$
     \ENDFOR
     \STATE Apply a guided DDIM step to obtain $a_{t-1}$
 \ENDFOR
 \STATE \textbf{return} $\mathrm{clip}(a_0,-1,1)$
 \end{algorithmic}
 }

\vspace{0.2em}
The corrector enforces the intended level-wise distribution; the predictor moves efficiently to the next cleaner diffusion level.
\end{block}

\end{column}

\begin{column}{\colwidth}

\begin{block}{\Large Results: MuJoCo Online RL}
\begin{center}
\includegraphics[width=0.98\textwidth]{eta45_vs_baselines.png}
\end{center}

\small
\vspace{0.2em}
\textbf{Setup:} six MuJoCo environments, 1M environment steps, multiple seeds, and comparisons against SAC, TD3, PPO, TRPO, DIPO, and the earlier diffusion-policy mirror-descent baseline DPMD.

 \vspace{0.3em}
 \textbf{Reading the figure:}
 MGMD is competitive across the benchmark suite and improves clearly over the earlier diffusion-policy mirror-descent approach, showing that inference-time MALA corrections can translate into better online control performance.
 \end{block}
 
 \vspace{0.8cm}
 
 \begin{block}{\Large Online RL Loop in the Current Codebase}
 \small
 The training loop follows the authoritative presentation closely:
 \begin{enumerate}
      \item Sample an action from the current tilted policy with the MALA-guided sampler.
      \item Execute it, collect $(s,a,r,s',d)$, and append to the replay buffer.
      \item Update twin critics by TD learning using \textbf{min}-$Q$ targets and a Huber loss.
      \item Fit the diffusion policy to sampled actions with the usual denoising loss.
 \end{enumerate}
 
 \vspace{0.3em}
 \textbf{Important implementation detail:}
 \begin{itemize}
     \item use the \textbf{mean} of the twin critics during sampling for a stronger guidance signal,
     \item use the \textbf{min} of the twin critics inside TD targets to limit overestimation.
 \end{itemize}
 
 \vspace{0.3em}
 \textbf{Representative settings used in this code:}
 \begin{itemize}
     \item 20 diffusion levels with a cosine schedule,
     \item deterministic DDIM predictor,
     \item 2 MALA steps per diffusion level,
     \item per-level Robbins--Monro adaptation of $s_t$,
     \item Tweedie/Q clip radius $3.0$, Huber width $30$, replay size $4\times 10^5$, UTD ratio $8$.
 \end{itemize}
 \end{block}
 
 \vspace{0.8cm}
 
 \begin{alertblock}{\Large Takeaways}
 \small
\begin{itemize}
    \item \textbf{Mirror descent target:} the desired policy update is an exponential tilt of the current diffusion policy by the critic.
    \item \textbf{Core technical move:} rewrite guided diffusion as a sequence of energy-based targets and correct each level with MALA.
    \item \textbf{Why it matters:} energy parameterization preserves standard diffusion training while making exact MH corrections possible.
    \item \textbf{Practical knob:} more inference-time compute can be spent on extra MALA steps, rather than changing the training objective.
\end{itemize}
\end{alertblock}

\end{column}

\end{columns}

\end{frame}
\end{document}
